\chapter{Euclid幾何学}
    \section{諸公式}
        \subsection{ヘロンの公式}
            \begin{shadebox}
                辺の長さが$a,b,c$なる三角形の面積$S$は
                \[S = \sqrt{s(s-a)(s-b)(s-c)}\]
                但し
                \[s = \frac{a+b+c}{2}\]
            \end{shadebox}
            \begin{proof}
                余弦定理を用いて
                \[\cos{C} = \frac{a^2+b^2-c^2}{2ab}\]
                \[\therefore \; \sin{C} = \sqrt{1-\cos^2{C}} = \frac{\sqrt{4a^2b^2-{\left(a^2+b^2-c^2\right)}^2}}{2ab}\]
                \begin{align*}
                    S &=& \frac{1}{2}ab\sin{C}\\
                    &=& \frac{1}{2}\sqrt{4a^2b^2-{\left(a^2+b^2-c^2\right)}^2}
                \end{align*}
                ここで
                \begin{align*}
                    &\;& 4a^2b^2-{\left(a^2+b^2-c^2\right)}^2\\
                    &=& -\left(a^4+b^4+c^4-2a^2b^2-2b^2c^2-2c^2a^2\right)\\
                    &=& -\left[a^4 - 2(b^2+c^2)a^2 + b^4+c^4 -2b^2c^2\right]\\
                    &=& -\left[a^4 - 2(b^2+c^2)a^2 + (b^2+c^2)^2 -4b^2c^2\right]\\
                    &=& -\left\{\left[a^2-(b^2+c^2)^2\right]^2-(2bc)^2\right\}\\
                    &=& -\left[a^2-(b^2+c^2)+2bc\right]\left[a^2-(b^2+c^2)-2bc\right]\\
                    &=& -\left[a^2-(b-c)^2\right]\left[a^2-(b+c)^2\right]\\
                    &=& (a+b+c)(-a+b+c)(a-b+c)(a+b-c)
                \end{align*}
                よって
                \begin{align*}
                    S &=& \sqrt{\frac{a+b+c}{2}\times\frac{-a+b+c}{2}\times\frac{a-b+c}{2}\times\frac{a+b-c}{2}}\\
                    &=& \sqrt{s(s-a)(s-b)(s-c)}
                \end{align*}
                但し
                \[s = \frac{a+b+c}{2}\]
            \end{proof}
    \section{諸定理}
        \subsection{\texorpdfstring{$\lim_{\|\bm{x}\|\to\infty}\|\bm{x}-\bm{a}\| - \|\bm{x}\| = -\frac{\bm{x}}{\|\bm{x}\|}\cdot\bm{a}$}{近接する2点と遠方の点との距離の差の極限}}
            \begin{shadebox}
                $\bm{a},\bm{x}\in\realNumbers^3$とするとき次式が成り立つ。
                \[ \lim_{\|\bm{x}\|\to\infty}\|\bm{x}-\bm{a}\| - \|\bm{x}\| = -\frac{\bm{x}}{\|\bm{x}\|}\cdot\bm{a} \]
            \end{shadebox}
            \begin{proof}
                \quad\par
                $\bm{a}=\bm{0}$のときは明らかに成り立つ。
                以下では$\bm{a}\neq 0$とする。
                また、$\|\bm{x}\|$を十分大きくとり、$\|\bm{x}\|>\|\bm{a}\|/2$とする。
                $f(\bm{x}) := \|\bm{x}-\bm{a}\| - \|\bm{x}\|$とすると次式が成り立つ。
                \[ \bigl(f(\bm{x}) + \|\bm{x}\|\bigr)^2 = \|\bm{x}-\bm{a}\|^2 \quad \therefore f(\bm{x})\bigl(f(\bm{x}) + 2\|\bm{x}\|\bigr) = \|\bm{a}\|^2 - 2\bm{x}\cdot\bm{a} \]
                $f(\bm{x})$が最小となるのは$\bm{a}$が、原点と$\bm{x}$を結ぶ線分の上にあるときで、その値は$-\|\bm{a}\|$である。
                $\|\bm{x}\|>\|\bm{a}\|/2$と仮定しているから次式を得る。
                \[ f(\bm{x}) = \frac{\|\bm{a}\|^2 - 2\bm{x}\cdot\bm{a}}{f(\bm{x}) + 2\|\bm{x}\|} < \frac{\|\bm{a}\|^2 - 2\bm{x}\cdot\bm{a}}{2\|\bm{x}\| - \|\bm{a}\|} \to -\frac{\bm{x}}{\|\bm{x}\|}\cdot\bm{a} \quad \text{as} \quad \|\bm{x}\|\to\infty\]
            \end{proof}